\chapter{O Operador Pipe}
\label{cap:pipes}
% ── índice ──────────────────────────────────────────────────────
\index{pipe, operador (	exttt{|>})|textbf}
\index{composição de funções}

\section{Sintaxe básica}

O operador \lstinline{|>} passa o valor à sua esquerda como primeiro argumento da
função à sua direita:

\begin{lstlisting}
expr |> f(args...)
// equivale a: f(expr, args...)
\end{lstlisting}

Sem pipe:
\begin{lstlisting}
display sort(filter(df, renda > 0), renda desc)
\end{lstlisting}

Com pipe:
\begin{lstlisting}
df |> filter(renda > 0) |> sort(renda desc) |> display
\end{lstlisting}

O pipe permite ler transformações na ordem em que ocorrem, da esquerda para a
direita (ou de cima para baixo), sem aninhamento de parênteses.

\section{Semântica com DataFrames}

O comportamento do pipe com DataFrames segue uma regra única:

\begin{center}
\begin{tabular}{ll}
\toprule
\textbf{Forma} & \textbf{Efeito} \\
\midrule
\texttt{df |> f(...)}          & Assign-back: df recebe o resultado \\
\texttt{let df2 = df |> f(...)} & Funcional: df inalterado, df2 é novo \\
\bottomrule
\end{tabular}
\end{center}

\subsection*{Assign-back}

Quando o pipe começa com uma variável (\lstinline{df |> ...}) e não há \lstinline{let},
o resultado é atribuído de volta à variável de origem:

\begin{lstlisting}[caption={Assign-back: df é modificado}]
load "dados.csv" as df

df |> filter(ano >= 2010)
   |> dropna
   |> sort(renda desc)

// df agora é o resultado de todas as transformações
\end{lstlisting}

\subsection*{Forma funcional com \texttt{let}}

\begin{lstlisting}[caption={let preserva o original}]
load "dados.csv" as df

let df_limpo = df
  |> filter(ano >= 2010)
  |> dropna

// df está intacto; df_limpo é a versão filtrada
\end{lstlisting}

\section{COW e pipes}

A semântica COW garante que o assign-back é eficiente: enquanto \lstinline{df} e
\lstinline{df_limpo} compartilham os mesmos dados (antes de qualquer mutação), nenhuma
cópia de memória ocorre. A cópia real acontece apenas quando os dados divergem.

\begin{lstlisting}[caption={Preservando baseline com let}]
load "painel.csv" as df
let baseline = df      // alias barato

df |> filter(tratamento == 1)   // modifica df, não baseline

display count(baseline)    // todas as obs originais
display count(df)          // só o grupo de tratamento
\end{lstlisting}

\section{Encadeamento longo}

O pipe é especialmente útil em pipelines de limpeza e transformação:

\begin{lstlisting}[caption={Pipeline completo de preparação}]
load "pnad.csv" as pnad

let amostra = pnad
  |> filter(idade >= 18 and idade <= 65)
  |> filter(horas_trab > 0)
  |> dropna
  |> generate(lrenda = log(renda + 1))
  |> generate(exp = idade - anos_estudo - 6)
  |> generate(exp2 = exp^2)
  |> sort(uf ano)

display count(amostra)
\end{lstlisting}

\section{Pipe com closures}

O operador \lstinline{|>} também funciona com closures:

\begin{lstlisting}
let resultado = 5 |> fn(x) { x^2 + 1 }
display resultado    // 26
\end{lstlisting}
